\documentclass[a4paper,12pt]{exam}
\usepackage[T1]{fontenc}
\usepackage{lmodern}
\usepackage{comment}
\usepackage{fancyvrb}
\usepackage{graphicx}
\usepackage{geometry}
\usepackage{amssymb,amsmath}
\usepackage{ifxetex,ifluatex}
\usepackage{fixltx2e} % provides \textsubscript
% use upquote if available, for straight quotes in verbatim environments
\IfFileExists{upquote.sty}{\usepackage{upquote}}{}
\ifnum 0\ifxetex 1\fi\ifluatex 1\fi=0 % if pdftex
  \usepackage[utf8]{inputenc}
\else % if luatex or xelatex
  \ifxetex
    \usepackage{mathspec}
    \usepackage{xltxtra,xunicode}
  \else
    \usepackage{fontspec}
  \fi
  \defaultfontfeatures{Mapping=tex-text,Scale=MatchLowercase}
  \newcommand{\euro}{€}
\fi
% use microtype if available
\IfFileExists{microtype.sty}{\usepackage{microtype}}{}
\usepackage{longtable,booktabs}
\usepackage{graphicx}
% Redefine \includegraphics so that, unless explicit options are
% given, the image width will not exceed the width of the page.
% Images get their normal width if they fit onto the page, but
% are scaled down if they would overflow the margins.
\makeatletter
\def\ScaleIfNeeded{%
  \ifdim\Gin@nat@width>\linewidth
    \linewidth
  \else
    \Gin@nat@width
  \fi
}
\makeatother
\let\Oldincludegraphics\includegraphics
{%
 \catcode`\@=11\relax%
 \gdef\includegraphics{\@ifnextchar[{\Oldincludegraphics}{\Oldincludegraphics[width=\ScaleIfNeeded]}}%
}%
\ifxetex
  \usepackage[setpagesize=false, % page size defined by xetex
              unicode=false, % unicode breaks when used with xetex
              xetex]{hyperref}
\else
  \usepackage[unicode=true]{hyperref}
\fi
\hypersetup{breaklinks=true,
            bookmarks=true,
            pdfauthor={},
            pdftitle={Worksheet: Streams},
            colorlinks=true,
            citecolor=blue,
            urlcolor=blue,
            linkcolor=magenta,
            pdfborder={0 0 0}}
\urlstyle{same}  % don't use monospace font for urls
\setlength{\parindent}{0pt}
\setlength{\parskip}{6pt plus 2pt minus 1pt}
\setlength{\emergencystretch}{3em}  % prevent overfull lines
%\setcounter{secnumdepth}{0}

\geometry{
  %includeheadfoot,
  margin=2.54cm
}

\title{Supplemental Exercises --- Week Five}
\date{Spring term 2017}
\author{}

\begin{document}
\maketitle

These exercises \textbf{are not} part of the coursework portfolio and are provided 
solely to give you further experience with the following topics:
\begin{itemize}
\item Pattern Matching, 
\item Class Arguments, 
\item Named \& Default Arguments,
\item Constructors, 
\item Auxiliary Constructors, 
\item Case Classes,
\item Parameterised Types, 
\item Functions as Objects, 
\item map \& reduce	
\end{itemize}
The source code for these exercises can be found under \verb!exercises/week05! together 
with some appropriate (minimalist) tests. Feel free to augment the tests as you feel fit.

\begin{questions}

\section*{Pattern Matching}	

\question
Create a method \texttt{forecast} that represents the percentage of cloudiness, 
and use it to produce a \verb!“weather forecast”! string such as \verb!“Sunny”! (100), 
\verb!“Mostly Sunny”! (80), \verb!“Partly Sunny”! (50), \verb!“Mostly Cloudy”! (20), 
and \verb!“Cloudy”! (0). For this exercise, only match for the legal 
values 100, 80, 50, 20, and 0. 
All other values should produce \verb!“Unknown”!. 

Satisfy the following tests:
\VerbatimInput{src/test/scala/forecast/ForecastTest.scala}

\begin{solution}
\VerbatimInput{src/main/scala/solutions/Forecast.scala}
\end{solution}

\question
Create a \texttt{Vector} named \texttt{sunnyData} that holds the values 
(100, 80, 50, 20, 0, 15). Use a \texttt{for} loop to call \texttt{forecast} 
with the contents of \texttt{sunnyData}. 
Display the answers and ensure that they match the responses above.
\begin{solution}
\VerbatimInput{src/main/scala/solutions/ForecastVector.scala}
\end{solution}

\section*{Class Arguments}

\question
Create a new class \texttt{Family} that takes a variable argument list 
representing the names of family members. 
Satisfy the following tests:
\VerbatimInput{src/test/scala/classargs/TestArgs.sc}

\begin{solution}
\VerbatimInput{src/main/scala/solutions/Family.scala}
\end{solution}

\question
Adapt the \texttt{Family} class definition to include class arguments for a mother, 
father, and a variable number of children. What changes did you have to make? 
Satisfy the following tests:
\begin{Verbatim}
val family3 = new FlexibleFamily("Mum", "Dad", "Sally", "Dick")
family3.familySize() is 4
val family4 = new FlexibleFamily("Dad", "Mom", "Harry") 
family4.familySize() is 3
\end{Verbatim}

\begin{solution}
\VerbatimInput{src/main/scala/solutions/FlexibleFamily.scala}
\end{solution}

\question
Write a method that squares a variable argument list of numbers and returns the sum. 
Satisfy the following tests:
\begin{Verbatim}
squareThem(2) is 4
squareThem(2, 4) is 20
squareThem(1, 2, 4) is 21
\end{Verbatim}

\begin{solution}
\VerbatimInput{src/main/scala/solutions/SquareThem.scala}
\end{solution}

\section*{Named \& Default Arguments}

\question
Define a class \texttt{SimpleTime} that takes two arguments: an \texttt{Int} 
that represents hours, and an \texttt{Int} that represents minutes. 
Use named arguments to create a \texttt{SimpleTime} object. Satisfy the following tests:
\begin{Verbatim}
val t = new SimpleTime(hours=5, minutes=30) 
t.hours is 5
t.minutes is 30
\end{Verbatim}

\begin{solution}
\begin{Verbatim}
class SimpleTime(val hours: Int, val minutes: Int)

val t = new SimpleTime(hours=5, minutes=30)	
\end{Verbatim}	
\end{solution}

\question
Using the solution for \texttt{SimpleTime} above, default \texttt{minutes} to 0 
so that you don’t have to specify them. Satisfy the following tests: 
\begin{Verbatim}
val t2 = new SimpleTime2(hours=10)
t2.hours is 10
t2.minutes is 0
\end{Verbatim}

\begin{solution}
\begin{Verbatim}
class SimpleTime2(val hours: Int, val minutes: Int = 0)

val t2 = new SimpleTime2(hours=10)	
\end{Verbatim}	
\end{solution}

\question
Create a class \texttt{Planet} that has, by default, a single moon. 
The \texttt{Planet} class should have a name (\texttt{String}) and description (\texttt{String}). 
Use named arguments to specify the name and description, and a default for the number of moons. 
Satisfy the following tests:
\begin{Verbatim}
val p = new Planet(name = "Mercury",
description = "small and hot planet", moons = 0)
p.hasMoon is false
\end{Verbatim}

\begin{solution}
\VerbatimInput{src/main/scala/solutions/PlanetTest.scala}
\end{solution}

\question
Modify your solution for the previous exercise by changing the order of the arguments 
that you use to create the \texttt{Planet}. Did you have to change any code? 
Satisfy the following tests:
\begin{Verbatim}
val earth = new Planet(moons = 1, name = "Earth",
    description = "a hospitable planet") 
earth.hasMoon is true
\end{Verbatim}

\begin{solution}
%\VerbatimInput{src/solutions/PlanetTestAgain.scala}
\end{solution}

\question
Demonstrate that named and default arguments can be used with methods. 
Create a class \texttt{Item} that takes two class arguments: 
A \texttt{String} for \texttt{name} and a \texttt{Double} for \texttt{price}. 

Add a method \texttt{cost} which has named arguments for \texttt{grocery (Boolean)}, 
\texttt{medication (Boolean)}, and \texttt{taxRate(Double)}. 
Default \texttt{grocery} and \texttt{medication} to \texttt{false}, 
\texttt{taxRate} to \texttt{0.10}. 

In this scenario, groceries and medications are not taxable. 
Return the total cost of the item by calculating the appropriate tax. 
Satisfy the following tests:
\begin{Verbatim}
val flour = new Item(name="flour", 4) 
flour.cost(grocery=true) is 4
val sunscreen = new Item(
name="sunscreen", 3)
sunscreen.cost() is 3.3
val tv = new Item(name="television", 500) 
tv.cost(rate = 0.06) is 530
\end{Verbatim}

\begin{solution}
\VerbatimInput{src/main/scala/solutions/ItemTest.scala}
\end{solution}

\section*{Constructors}

\question
Create a new class \texttt{Tea} that has two methods: 
\begin{description}
\item[\texttt{describe}] --- which includes information about whether the tea 
includes milk, sugar, is decaffeinated, and includes the name; and 
\item[\texttt{calories}] --- which adds 100 calories for milk and 16 calories for sugar. 
\end{description}
Satisfy the following tests:
\VerbatimInput{src/test/scala/tea/TeaScript.sc}

\begin{solution}
\VerbatimInput{src/main/scala/solutions/TeaScript.scala}
\end{solution}

\section*{Auxiliary Constructors}

\question
Create a class called \texttt{ClothesWasher} with a default constructor 
and two auxiliary constructors, one that specifies \texttt{modelName} 
(as a \texttt{String}) and one that specifies \texttt{capacity} (as a \texttt{Double}).

\begin{solution}
\VerbatimInput{src/main/scala/solutions/ClothesWasherTest.scala}
\end{solution}

\question
Create a class \texttt{ClothesWasher2} that looks just like your solution above, 
but use named and default arguments instead so that you can produce the same results with 
just a default constructor.

\begin{solution}
\VerbatimInput{src/main/scala/solutions/ClothesWasherTest2.scala}
\end{solution}

\question
Show that the first line of an auxiliary constructor must be a call to the primary constructor.

\begin{solution}
\VerbatimInput{src/main/scala/solutions/TryMeTest.scala}
\end{solution}

\section*{Case Classes}

\question
Create a \emph{case class} to represent a \texttt{Person} in an address book, 
complete with \texttt{String}s for the name and a \texttt{String} for contact information. 
Satisfy the following tests:
\begin{Verbatim}
val p = Person("Jane", "Smile", "jane@smile.com")
p.first is "Jane"
p.last is "Smile"
p.email is "jane@smile.com"
\end{Verbatim}

\begin{solution}
\begin{Verbatim}
case class Person(first: String, last: String, email: String)
val p = Person("Jane", "Smile", "jane@smile.com")	
\end{Verbatim}	
\end{solution}

\question
Create some \texttt{Person} objects. 
Put the \texttt{Person} objects in a \texttt{Vector}. 
Satisfy the following tests:
\begin{Verbatim}
val people = Vector( 
    Person("Jane","Smile","jane@smile.com"), 
    Person("Ron","House","ron@house.com"), 
    Person("Sally","Dove","sally@dove.com"))

people(0) is "Person(Jane,Smile,jane@smile.com)"
people(1) is "Person(Ron,House,ron@house.com)" 
people(2) is "Person(Sally,Dove,sally@dove.com)"
\end{Verbatim}

\begin{solution}
\VerbatimInput{src/main/scala/solutions/PersonTest.scala}
\end{solution}

\section*{Parameterised Types}

\question
Modify the method \texttt{explicit} in the following code:
\VerbatimInput{src/main/scala/parameter/ParameterizedReturnTypes.sc} 
so it creates and returns a \texttt{Vector} of \texttt{Double}. 
Satisfy the following tests: 
\begin{Verbatim}
explicitDouble(1.0, 2.0, 3.0) is Vector(1.0, 2.0, 3.0)
\end{Verbatim}

\begin{solution}
\begin{Verbatim}
def explicitDouble(
  v1:Double, v2:Double, v3:Double):Vector[Double] = {
  Vector(v1, v2, v3)
}	
\end{Verbatim}	
\end{solution}

\question
Building on the previous exercise, alter \texttt{explicit} to take a 
\texttt{Vector} and create and return a \texttt{List}. 
Refer to the \texttt{ScalaDoc} for \texttt{List}, if necessary. 

Satisfy the following tests: 
\begin{Verbatim}
explicitList(Vector(10.0, 20.0)) is List(10.0, 20.0)

explicitList(Vector(1, 2, 3)) is List(1.0, 2.0, 3.0)
\end{Verbatim}
\begin{solution}
\begin{Verbatim}
def explicitList(v1:Vector[Double]):List[Double] = {
    v1.toList
}	
\end{Verbatim}	
\end{solution}

\question
Building on the previous exercise, change \texttt{explicit} to return a \texttt{Set}. 

Satisfy the following tests:
\begin{Verbatim}
explicitSet(Vector(10.0, 20.0, 10.0)) is Set(10.0, 20.0)
explicitSet(Vector(1, 2, 3, 2, 3, 4)) is Set(1.0, 2.0, 3.0, 4.0)
\end{Verbatim}
\begin{solution}
\begin{Verbatim}
def explicitSet(v1:Vector[Double]):Set[Double] = {
    v1.toSet
}	
\end{Verbatim}	
\end{solution}

\section*{Functions as Objects}

%\question
%Working from your solution to the exercise above, add a comma between each number. 
%Satisfy the following test:
%\begin{Verbatim}
%str is "1,2,3,4,"
%\end{Verbatim}
%
%\begin{solution}
%\begin{Verbatim}
%var str1 = ""
%val numberV = Vector(1, 2, 3, 4)
%numberV.foreach(n=>str1 += n + ",")
%str1 is "1,2,3,4,"	
%\end{Verbatim}	
%\end{solution}

\question
Create an anonymous function that calculates age in \emph{dog years} 
(by multiplying years by 7). 
Assign it to a \verb!val! and then call your function. 

Satisfy the following test:
\begin{Verbatim}
val dogYears = // Your function here
dogYears(10) is 70
\end{Verbatim}

\begin{solution}
\begin{Verbatim}
val dogYears = (x:Int) => (x*7)
dogYears(10) is 70	
\end{Verbatim}	
\end{solution}

\question
Create a \texttt{Vector} and append the result of \texttt{dogYears} to a \texttt{String} 
for each value in the \texttt{Vector}. 

Satisfy the following test:
\begin{Verbatim}
var s = ""
val v = Vector(1, 5, 7, 8)
v.foreach(/* Fill this in */)
s is "7 35 49 56 "
\end{Verbatim}

\begin{solution}
\begin{Verbatim}
val v = Vector(1, 5, 7, 8)
val dogYears = (x:Int) => (x*7)

var s = ""
v.foreach(x=>(s = s + (dogYears(x) + " ")))
s is "7 35 49 56 "	
\end{Verbatim}	
\end{solution}

\question
Create an anonymous function to square a list of numbers. 
Call the function for every element in a \texttt{Vector}, using \texttt{foreach}. 

Satisfy the following test:
\begin{Verbatim}
var s = ""
val numbers = Vector(1, 2, 5, 3, 7) 
numbers.foreach(/* Fill this in */) 
s is "1 4 25 9 49 "
\end{Verbatim}

\begin{solution}
\begin{Verbatim}
var s = ""
val numbers = Vector(1, 2, 5, 3, 7)
numbers.foreach(x=>(s = s +(x*x + " ")))	
\end{Verbatim}	
\end{solution}

\end{questions}
\end{document}